%% it/sections/02_regole.tex
%% ============================================================
%% 02_regole.tex — Regole FitD adattate / Adapted FitD Rules
%% ============================================================

\chapter{Le Regole}
\markboth{Le Regole}{Le Regole}

\lettrine[lines=3]{V}{espri} 1282 usa il sistema \textbf{Forged in the Dark}
(FitD), creato da John Harper per \textit{Blades in the Dark}. Questa sezione
presenta le regole adattate al contesto storico e alla struttura multi-tavolo.

\section{Il Cuore del Sistema}

Quando un personaggio tenta qualcosa di rischioso e l'esito è incerto,
si tira un dado. Si raccoglie un numero di dadi a sei facce pari al valore
dell'azione pertinente, si tira, e si considera il risultato più alto.

\regola[Tiro di Azione]{
  \textbf{6}: Successo pieno. Ottieni quello che vuoi, senza costo.\\
  \textbf{4--5}: Successo parziale. Ottieni quello che vuoi, ma con
  una complicazione, un costo, o un effetto ridotto.\\
  \textbf{1--3}: Fallimento. Le cose vanno storte. Il GM introduce
  una conseguenza negativa.\\[4pt]
  Se non hai dadi (azione a 0), tira \textbf{due dadi} e prendi il
  risultato più basso.
}

\subsection{Posizione e Effetto}

Prima di ogni tiro, il GM stabilisce due parametri fondamentali:

\textbf{Posizione} — quanto sei esposto al rischio:
\begin{itemize}[nosep]
  \item \textit{Controllata} — sei al sicuro, il rischio è minimo
  \item \textit{Rischiosa} — situazione tesa, le conseguenze sono serie
  \item \textit{Disperata} — sei in pericolo reale, le conseguenze possono essere devastanti
\end{itemize}

\textbf{Effetto} — quanto riesci a incidere sulla situazione:
\begin{itemize}[nosep]
  \item \textit{Limitato} — il tuo impatto è piccolo
  \item \textit{Standard} — il tuo impatto è quello atteso
  \item \textit{Eccezionale} — il tuo impatto supera le aspettative
\end{itemize}

\nota{Per il GM}{
  Posizione ed effetto sono strumenti narrativi, non parametri fissi.
  Cambiali liberamente in risposta alle azioni dei giocatori. Un piano
  ben elaborato può spostare una posizione da Rischiosa a Controllata;
  una distrazione improvvisa può fare il contrario.
}

\secsep

\section{Le Azioni}

I personaggi di \textit{Vespri 1282} agiscono attraverso \textbf{dodici
azioni} raggruppate in quattro attributi che riflettono le tensioni
del periodo storico.

\subsection{Cuore — \textit{Passione e Volontà}}

\begin{description}[leftmargin=2em, style=nextline]
  \item[\cinzel Aizzare] Infiammare gli animi, muovere le folle,
    fare appelli emotivi. L'azione del predicatore, del demagogo, del martire.
  \item[\cinzel Resistere] Tenere duro sotto pressione, sopportare il
    dolore fisico o emotivo, non cedere all'interrogatorio.
  \item[\cinzel Sfidare] Affrontare direttamente un avversario, mettersi
    davanti al pericolo, accettare il confronto aperto.
\end{description}

\subsection{Mente — \textit{Intelligenza e Astuzia}}

\begin{description}[leftmargin=2em, style=nextline]
  \item[\cinzel Studiare] Raccogliere informazioni, leggere persone e
    situazioni, capire i piani altrui.
  \item[\cinzel Tramare] Elaborare piani, costruire reti, coordinare
    agenti su scala larga.
  \item[\cinzel Ingannare] Mentire con convinzione, costruire false
    identità, confondere gli avversari.
\end{description}

\subsection{Corpo — \textit{Forza e Destrezza}}

\begin{description}[leftmargin=2em, style=nextline]
  \item[\cinzel Agire] Muoversi rapidamente, eseguire azioni fisiche
    di precisione, furtività.
  \item[\cinzel Combattere] Usare la violenza in modo efficace, proteggere
    gli alleati, disarmare o neutralizzare avversari.
  \item[\cinzel Manovrare] Spostarsi attraverso spazi controllati,
    passare inosservati, sfuggire alla sorveglianza.
\end{description}

\subsection{Spirito — \textit{Influenza e Fede}}

\begin{description}[leftmargin=2em, style=nextline]
  \item[\cinzel Comandare] Guidare con autorità, dare ordini che vengono
    eseguiti, mantenere la disciplina.
  \item[\cinzel Persuadere] Convincere con argomenti, negoziare, sedurre con parole.
  \item[\cinzel Invocare] Appellarsi a forze più grandi — Dio, la giustizia,
    la memoria degli antenati. Azione speciale del Clero, disponibile
    a tutti nei momenti estremi.
\end{description}

\secsep

\section{Stress e Trauma}

Ogni personaggio ha \textbf{9 caselle di Stress}. Lo stress si accumula
quando si subiscono conseguenze, si usa la fortuna, ci si mette in
situazioni disperate. Quando si riempie la nona casella, il personaggio
subisce un \textbf{Trauma} e la casella si svuota.

\subsection{Traumi di Vespri 1282}

\begin{itemize}[nosep]
  \item \textit{Sfiducia} — non ti fidi più di nessuno, nemmeno degli alleati
  \item \textit{Ossessione} — un obiettivo consuma ogni altra considerazione
  \item \textit{Colpa} — hai fatto qualcosa di imperdonabile
  \item \textit{Paura} — un tipo specifico di situazione ti paralizza
  \item \textit{Fede perduta} — non credi più in ciò che ti ha portato fin qui
\end{itemize}

Un personaggio con \textbf{4 traumi} è fuori gioco.

\regola[Usare la Fortuna]{
  Prima di un tiro, puoi dichiarare di usare la Fortuna. Aggiungi
  \textbf{1d6} al tiro. Poi subisci \textbf{2 punti di Stress}.
  Puoi usare la Fortuna più volte sullo stesso tiro (massimo 3 volte).
}

\regola[Resistere alle Conseguenze]{
  Tira i dadi dell'attributo pertinente. Ogni risultato sopra 3 riduce
  la severità della conseguenza. Un 6 la elimina completamente. Subisci
  comunque \textbf{1 punto di Stress} per aver resistito.
}

\secsep

\section{Meccaniche Multi-Tavolo}

\subsection{Il Punteggio di Fazione}

Ogni fazione accumula un punteggio da 0 a 4 durante la sessione.
Il punteggio totale dei quattro tavoli (0--16) determina l'esito
finale della notte del Vespro.

\nota{Per il GM Coordinatore}{
  Il GM coordinatore tiene traccia dei punteggi di tutti i tavoli su una
  scheda centrale. I punteggi non vengono rivelati ai giocatori durante
  la sessione: sono la misura di ciò che sta accadendo fuori scena.
  Ogni tavolo sa solo quello che il suo GM comunica narrativamente.
}

\subsection{Momenti di Congiunzione}

In tre momenti della sessione il GM coordinatore comunica a tutti i tavoli
un \textbf{Evento Condiviso}: una notizia o conseguenza delle azioni
degli altri tavoli che entra nella narrativa di ciascun gruppo.

\begin{enumerate}[nosep]
  \item \textbf{La Voce} (dopo 45 min) — una notizia che cambia il quadro
  \item \textbf{La Crisi} (dopo 90 min) — qualcosa va storto su scala cittadina
  \item \textbf{Le Campane} (dopo 150 min) — la processione di Santo Spirito
    sta per cominciare
\end{enumerate}

\subsection{Il Bot Telegram}

\nota[Comandi Bot]{
  \texttt{/punteggio [fazione] [0-4]} — aggiorna punteggio fazione\\
  \texttt{/evento [testo]} — invia evento condiviso a tutti i tavoli\\
  \texttt{/totale} — mostra punteggio aggregato (solo coordinatore)\\
  \texttt{/stato} — mostra stato corrente di tutti i tavoli
}

\secsep

\section{Componenti Fisici Opzionali}

\textit{Vespri 1282} si gioca benissimo con carta e matita. Ma per una
sessione dal vivo — soprattutto un evento serale a più tavoli — alcuni
componenti fisici rendono le meccaniche tangibili. Sono del tutto
opzionali e non cambiano le regole: rendono solo visibile ciò che accade.

\nota[Props consigliati]{
  \textbf{Gettoni Stress} — nove pietre scure per giocatore, da posare
  sulle caselle Stress della scheda. La scheda resta il riferimento;
  i gettoni rendono visibile il peso che il personaggio porta.\\[4pt]
  \textbf{Gettoni Passione} — cinque perline rosse per giocatore, per
  chi usa il layer Passionale (vedi «La Passione dei Vespri»). Spendere
  Fortuna, resistere o entrare in Crisi diventano gesti visibili al tavolo.\\[4pt]
  \textbf{Nastri delle Connessioni} — un nastro o cordoncino per ogni
  Connessione Passionale dichiarata. Un nastro teso fra due tavoli
  segnala a tutta la sala che qualcosa li lega.
}

Bastano due colori di gettoni: la Fortuna spende Stress, non serve un
terzo tipo. Tieni le forbici per i nastri alla postazione del
coordinatore, non ai tavoli: recidere un legame dev'essere un momento
deliberato.

\filettodoppio
